\documentclass[12pt,a4paper]{article}
\usepackage[utf8]{inputenc}
\usepackage[spanish]{babel}
\usepackage[margin=2.5cm]{geometry}
\usepackage{hyperref}

\hypersetup{
    colorlinks=true,
    urlcolor=blue
}

\pagestyle{empty}

\begin{document}

\begin{center}
    {\Large \textbf{UNIVERSIDAD TÉCNICA ESTATAL DE QUEVEDO}}\\[0.5em]
    {\large FACULTAD DE CIENCIAS DE LA COMPUTACIÓN}\\[0.3em]
    {\large CARRERA DE SOFTWARE}\\[1.5em]
    
    {\large \textbf{INGENIERÍA DE REQUERIMIENTOS (ISR-401)}}\\[0.5em]
    {\normalsize 4.\textsuperscript{o} Nivel de Software --- Paralelo A}\\[2em]
    
    {\Large \textbf{HOJA DE ENTREGA DE ACTIVIDAD FORMATIVA}}\\[0.8em]
    {\large \textbf{FORM-TL-01: Taller de tipos de requisitos}}\\[2.5em]
\end{center}

\vspace{1cm}

\noindent \textbf{Estudiante:} Luis Alberto Paca Pilataza\\[1em]
\noindent \textbf{Docente:} Gleiston Cicerón Guerrero Ulloa, PhD\\[1em]
\noindent \textbf{Código y nombre de la tarea:} FORM-TL-01 --- Taller de tipos de requisitos\\[1em]
\noindent \textbf{Fecha límite de entrega:} Sábado 31 de octubre de 2026, a las 23:59\\[1.5em]
\noindent \textbf{URL del repositorio público en GitHub (en una sola línea):}\\
\url{https://github.com/Luispaca2002/ISR401-PacaLuis}

\vfill

\begin{center}
    \footnotesize Quevedo, Los Ríos, Ecuador --- Periodo Académico 2026--2027
\end{center}

\end{document}
